
%-------------------------------------------------------------------------------
\chapter{Data and Monte Carlo Samples}
\label{sec:samples}
%-------------------------------------------------------------------------------

This analysis is done using dedicated data samples collected by the ATLAS experiment in low pile-up conditions. Many of the data and MC samples are similar to those used for the W \pT measurement~\cite{STDM-2018-17} and the previous setup was further described in this support note~\cite{Kretzschmar:2657141}. As there were several smaller updates, this section summarises all key information and gives the latest status. 

\section{Data Sample}

In November 2017 ATLAS collected low pile-up data at
$\sqrt{s} = 5.02$~\TeV and $\sqrt{s} = 13$~\TeV, corresponding to data
periods M and N respectively.\footnote{The data taken at
  $\sqrt{s} = 5.02$~\TeV\ are often referred to as ``5~\TeV\ data''.}
In July 2018, ATLAS recorded another low pile-up dataset at
$\sqrt{s} = 13$~\TeV. The $13$~\TeV\ samples were levelled at an average
number of interaction per crossing, $\langle\mu\rangle \approx 2$,
while the $5$~\TeV\ data featured a similar average value with
levelling as well. See Ref~\cite{Kretzschmar:2657141} for the $\mu$
and the number of reconstructed primary vertices $N_\mathrm{PV}$
distributions. All data is taken at 25ns bunch spacing, but there is a
small difference in the filling scheme where the 2017 data was taken
in the $8b4e$ bunch structure and the 2018 data with longer
uninterrupted bunch trains.

\begin{itemize}
\item $\sqrt{s} = 5.02\tev$ data taken in November 2017, ATLAS data
  period M, final calibrated luminosity $254.938\,\ipb$ with an uncertainty
  of $\pm 1.0\%$ (tag
  OflLumi-5TeV-004), \\GRL
  \href{http://atlasdqm.web.cern.ch/atlasdqm/grlgen/All_Good/data17_5TeV.periodAllYear_DetStatus-v98-pro21-16_Unknown_PHYS_StandardGRL_All_Good_25ns_ignore_GLOBAL_LOWMU.xml}{\scriptsize data17\_5TeV.periodAllYear\_DetStatus-v98-pro21-16\_Unknown\_PHYS\_StandardGRL\_All\_Good\_25ns\_ignore\_GLOBAL\_LOWMU.xml}:\\
  Runs 340644, 340683, 340697, 340718, 340814, 340849, 340850, 340910, 340925, 340973, 341027, 341123, 341184
\item $\sqrt{s} = 13\tev$ data taken in November 2017, ATLAS data
  period N, final calibrated luminosity $147.061\ipb$ with an uncertainty
  of $\pm 1.15\%$ (tag
  OflLumi-13TeV-011), \\GRL
  \href{http://atlasdqm.web.cern.ch/atlasdqm/grlgen/All_Good/data17_13TeV.periodN_DetStatus-v98-pro21-16_Unknown_PHYS_StandardGRL_All_Good_25ns_ignore_GLOBAL_LOWMU.xml}{\scriptsize data17\_13TeV.periodN\_DetStatus-v98-pro21-16\_Unknown\_PHYS\_StandardGRL\_All\_Good\_25ns\_ignore\_GLOBAL\_LOWMU.xml}:\\
  Runs 341294, 341312, 341419, 341534, 341615, 341649
\item $\sqrt{s} = 13\tev$ data taken in June 2018, ATLAS data
  period G4+J, final calibrated luminosity $ 191.085\ipb$ with an uncertainty
  of $\pm 1.08\%$ (tag
  OflLumi-13TeV-011), \\GRL
  \href{http://atlasdqm.web.cern.ch/atlasdqm/grlgen/All_Good/data18_13TeV.periodG4J_MERGED_PHYS_StandardGRL_All_Good_25ns_ignore_GLOBAL_LOWMU.xml}{\scriptsize data18\_13TeV.periodG4J\_MERGED\_PHYS\_StandardGRL\_All\_Good\_25ns\_ignore\_GLOBAL\_LOWMU.xml}:\\
  Runs 354396, 355331, 355389, 355416, 355468
\end{itemize}

The total luminosity for the 13 TeV data amounts to $338.146\ipb$ with
an uncertainty of $\pm 0.92\%$. The luminosity calibration of the 13
\TeV low-pileup datasets is discussed in Ref.~\cite{DAPR-2021-01} and
the calibration of the $5.02$ \TeV data sample was carried out in a
similar way.

There are two places in the reconstruction of signals in the ATLAS
calorimeters that are sensitive to the level of pile-up. During data
taking, the so-called LAr ``optimal filtering coefficients'',
\texttt{OFC}, and the settings of the readout gain for the EMEC-IW and
FCAL calorimeters front-end electronics~\cite{LARG-2009-01} affect the
``online noise-suppression''.  For the nominal Run 2 data and MC16
simulations, the LAr \texttt{OFC} were optimised for
$\langle\mu\rangle = 20$, the EMEC-IW and FCAL front-end electronics
are allowed to use only Medium-Gain and Low-Gain readouts, and \sigPU\
was taken to correspond to an average of $40$ additional $pp$
interactions per bunch crossing. For the \lowmu\ data used in this
measurement, the setting were kept identical, as studies showed the
impact was marginal and it would likely be beneficial to keep the
setting used throughout Run 2. However, in the offline reconstruction
of calorimeter clusters, the value of \sigPU\ is chosen to correspond
to $\langle\mu\rangle = 0$, \textit{i.e.} $\sigPU = 0$ to retain the
calorimeter sensitivity to the small signals relevant to the hadronic
recoil reconstruction.


The data of the 2017 runs, both 5 and 13 \TeV, was reprocessed using
updated settings for the ID alignment (\texttt{r10803}) to improve the
muon sagitta biases, while the 2018 data relies on the
T0-reconstructed data sets. Data is further processed to STDM4 or
STDM6 derivations:
\begin{itemize}
\item original (T0) 2017 data was processed to STDM4 with \texttt{p3372}
\item reprocessed 2017 data was processed to STDM6 with \texttt{p3653}
\item 2018 data was processed to STDM6 with \texttt{p3583}
\end{itemize}

\section{MC Samples and cross sections}
\label{sec:mcsamples}

A dedicated MC16 campaign was produced to match the special condition
of the \lowmu data.  All signal and background samples where
reprocessed using the following tags: \texttt{r10243} for 5~\TeV,
\texttt{r10244} for 13~\TeV\ 2017 data, and \texttt{r11165} for
13~\TeV\ 2018 data. 

All MC samples uses settings specific to the special data run conditions, 
i.e. specifically the pileup overlay, topo-cluster noise thresholds and trigger 
menu are adapted. Given that the pileup distribution is already adjusted closely to the dataset, 
no further $\mu$-reweighting is performed.

For $\sqrt{s} = 13\tev$, generally the event samples generated and
simulated (``HITS'', tag \texttt{s3126}) for the main ATLAS MC16d/e campaign are reused and
processed further with dedicated settings in digitisation and
reconstruction. For $\sqrt{s} = 5\tev$, event samples are processed specifically for
this data set, using the simulation and reconstruction tags
\texttt{s3238\_r10243}.

Tables~\ref{tab:samples13} and ~\ref{tab:samples5} summarize the
main simulated samples and their properties. The cross-sections quoted in
the table are used to normalize predicted event counts for the various
samples and the respective uncertainties are listed as well.

The main signal event samples for $W$ and $Z$ production are generated
at NLO QCD and LO EW using the \Powheg\ event
generator~\cite{Nason:2004rx,Frixione:2007vw,Alioli:2010xd,Barze:2012tt,Barze:2013fru}
using the CT10 PDF, interfaced to \Pythia~\cite{pythia} using the
AZNLO tune~\cite{STDM-2012-23}. \Powheg+\Pythia samples are interfaced
to \Photos~\cite{Golonka:2005pn} to simulate the effect of final state
QED radiation using NLO kernels and including the effect of low-mass
$\gamma^* \to \ell\ell$ splitting. In contrast to earlier such samples
produced with the older \Powheg~V1 code~\cite{Alioli:2008gx}, these
newer samples improve make use of the \Powheg~V2 modules that allow
PDF variations to be incorporated on the fly. Also, the \Photos{}
accuracy is formally improved.

Alternative signal event samples for $W$ and $Z$ production are
generated at NNLO QCD and LO EW using the MiNNLOPS
procedure~\cite{Monni:2020nks,Monni:2019whf} using the CT18NNLO
PDF~\cite{Hou:2019efy}, interfaced to \Pythia~\cite{pythia} using a
variation of the A14 tune~\cite{ATL-PHYS-PUB-2014-021}~\footnote{The
  parameters changed with respect to the ``vanilla'' A14 settings are
  Pythia's primordial $k_t$ parameter and the parameters affecting the
  parton shower matching. These parameters were chosen to optimise the improvement
  with existing unfolded $Z$ \pt and $Z+$jets data.}.
\Photos{} is used in the same mode as for the above  \Powheg+\Pythia samples.

Alternative samples at $\sqrt{s} = 13\tev$ are prepared with \Sherpa\
2.2.2/5~\cite{Hoche:2010kg} using the NNPDF3.0 PDFs and merging
$V+0,1,2$ at NLO accuracy with $V+3,4$ at LO accuracy using the
MEPS@NLO scheme. Futher technical details and overview talbles on
these samples are given in Ref.~\cite{Kretzschmar:2657141}.

Backgrounds from top-quark pair-production $t\bar{t}$ as well as
single-top production ($Wt$, t-channel, s-channel) are generated with
\Powheg+\Pythia. Backgrounds from di-bosons
$WW, WZ, ZZ$ are generated with \Powheg+\Pythia~\cite{Nason:2013ydw}
in all decay channels with at least one charged lepton in the final
state using the NNPDF3.0 NLO PDF set and the A14 Pythia tune. Special
care is taken to avoid overlaps in final states with the $W$ and $Z$
signal samples where an additional charged lepton pair is generated by
$\gamma^* \to \ell\ell$ splitting.

For muon calibration purposes, some dedicated $J/\psi \to \mu\mu$
samples are produced with \Pythia, interfacted to \Photos{} to
simulate QED FSR, both for the ``direct'' production process as well
as for ``indirect'' production in decays of B hadrons.

Cross sections for all processes listed in the tables below are
straight from the generator, except for the $t\bar{t}$ process, where
improved NNLO predictions are used directly. Specifically for the $W$
and $Z$ processes, corrections are applied later on as part of the
physics modelling procedures. Uncertainties on the cross-sections
arise from the choice of PDF set and PDF set internal uncertainties,
from factorization and renormalisation scale dependence, and the
strong coupling constant $\alpha_s$. Total uncertainties of $5\%$ ($W$
and $Z$), 7\% ($t\bar{t}$) and 10\% (single top, dibosons) are taken
on predictions normalized using these cross-sections,
e.g. for use of estimation of backgrounds.

The effect of multiple interactions per bunch crossing (``pileup'') is
modelled by overlaying simulated minimum bias events over the original
hard-scattering event. The minimum bias events were generated with
\Pythia\ using the NNPDF2.3LO set of PDFs~\cite{Ball:2012cx} and the A3
tune~\cite{ATL-PHYS-PUB-2016-017}.


\begin{table}[htbp]
  \small
  \begin{center}
    \begin{tabular}{l|l|l|l|l}
      \hline
      Process & Data set & $\sigma{\cdot}
      \text{BR}{\cdot}\epsilon_\mathrm{filter}$ [nb] (th. unc.)
      & $N^\mathrm{skim}_\mathrm{evt}\,[10^6]$
      & $N^\mathrm{unskim}_\mathrm{evt}\,[10^6]$\\
      \hline
      \multicolumn{5}{l}{\textbf{\Powheg+\Pythia NLO $W,Z$ samples}}\\
      \hline
$ W^{+} \to e^{+}\nu $ & 603391 & 11.39 (5\%)  & 45+58 & 45+58 \\
$ W^{+} \to \mu^{+}\nu $ & 603392 & 11.39 (5\%)  & 45+58 & 40+58 \\
$ W^{+} \to \tau^{+}\nu $ & 603393 & 11.39 (5\%)  & 20+26 & 20+26 \\
$ W^{-} \to e^{-}\bar{\nu} $ & 603394 & 8.341 (5\%)  & 35+44 & 35+44 \\
$ W^{-} \to \mu^{-}\bar{\nu} $ & 603395 & 8.341 (5\%)  & 35+45 & 35+45 \\
$ W^{-} \to \tau^{-}\bar{\nu} $ & 603396 & 8.341 (5\%)  & 15+13 & 15+13 \\\hline
$ Z \to ee $ $m>60\GeV$ & 603397 & 1.903 (5\%)  & 15+19 & 15+19 \\
$ Z \to \mu\mu $ $m>60\GeV$ & 603398 & 1.903 (5\%)  & 15+19  & 15+19 \\
$ Z \to \tau\tau $ $m>60\GeV$ & 603399 & 1.903 (5\%)  & 3.5+4.4 & 3.5+4.4 \\
$ Z \to ee $ $20<m<60\GeV$ & 603400 & 1.206 (5\%)  & 2.0+2.6 & 2.0+2.6 \\
$ Z \to \mu\mu $ $20<m<60\GeV$ & 603401 & 1.206 (5\%)  & 2.0+2.6 & 2.0+2.6 \\
      \hline
      \multicolumn{5}{l}{\textbf{\Powheg+\Pythia MiNNLOPS $W,Z$ samples -- Minitrees stil TODO}}\\
      \hline
$ W^{+} \to e^{+}\nu $ & 604351 & 11.64 (5\%) &  & 18.5+[TODO] \\
$ W^{+} \to \mu^{+}\nu $ & 604352 & 11.64  (5\%) &  & 18.4+23.4 \\
$ W^{+} \to \tau^{+}\nu $ & 604353 & 11.64  (5\%) &  & 8.4+10.6 \\
$ W^{-} \to e^{-}\bar{\nu} $ & 604348 & 8.603 (5\%) & & 14.5+18.5 \\
$ W^{-} \to \mu^{-}\bar{\nu} $ & 604349 & 8.603  (5\%) &  & 14.5+18.5 \\
$ W^{-} \to \tau^{-}\bar{\nu} $ & 604350 & 8.603  (5\%) &  & 6.3+7.9 \\\hline
$ Z \to ee $ $m>60\GeV$ & 604354 & 1.912 (5\%) &  & 6.1+7.8 \\
$ Z \to \mu\mu $ $m>60\GeV$ & 604355 & 1.912 (5\%) &  & 6.2+7.8 \\
$ Z \to \tau\tau $ $m>60\GeV$ & 604356 & 1.912 (5\%) &  & 1.4+1.8 \\
      \hline
      \multicolumn{5}{l}{\textbf{\Powheg+\Pythia Diboson samples}}\\
      \hline
$ WW (lvlv) $ & 604493 & 0.01055 (10\%)  & 0.037+0.037 & 0.05+0.05 \\
$ WW (lvqq) $ & 604494 & 0.04383 (10\%)  & 0.100+0.100 & 0.2+0.2 \\
$ WZ (lvll) $ & 604495 & 0.00144 (10\%)  & 0.008+0.008 & 0.01+0.01 \\
$ WZ (lvqq) $ & 604496 & 0.01021 (10\%)  & 0.026+0.026 & 0.05+0.05 \\
$ WZ (lvvv) $ & 604497 & 0.00280 (10\%)  & 0.010+0.010 & 0.02+0.02 \\
$ WZ (qqll) $ & 604498 & 0.00702 (10\%)  & 0.011+0.011 & 0.03+0.03 \\
$ ZZ (llll) $ & 604499 & 0.00061 (10\%)  & 0.004+0.004 & 0.01+0.01 \\
$ ZZ (qqll) $ & 604500 & 0.00609 (10\%)  & 0.011+0.011 & 0.03+0.03 \\
$ ZZ (vvll) $ & 604501 & 0.00093 (10\%)  & 0.005+0.005 & 0.01+0.01 \\
      \hline
      \multicolumn{5}{l}{\textbf{\Powheg+\Pythia top samples}}\\
      \hline
$ t\bar{t} $ & 410470 & 0.8318 $\times$ 0.544 (7\%)  & 1.35+1.28 & 2.0+1.9 \\
$ t $ (s-chan.) & 410644 & 0.002027 (10\%)  & 0.011+0.0055 & 0.02+0.01 \\
$ \bar{t} $ (s-chan.) & 410645 & 0.001268 (10\%)  & 0.0056+0.0056 & 0.01+0.01 \\
$ t $ (t-chan.) & 410658 & 0.03699 (10\%)  & 0.113+0.113 & 0.2+0.2 \\
$ \bar{t} $ (t-chan.) & 410659 & 0.02217 (10\%)  & 0.088+0.088 & 0.15+0.15 \\
$ W\bar{t} $ & 601352 & 0.03597 (10\%)  & 0.088+0.083 & 0.2+0.2 \\
$ Wt $ & 601355 & 0.03600 (10\%)  & 0.082+0.083 & 0.2+0.2\\
      \hline
      \multicolumn{5}{l}{\textbf{\Pythia $J/\psi \to \mu\mu$ samples}}\\
      \hline
Prompt & 801912 & -- & 10+13 & 10+13 \\
Non-prompt & 801917 & -- & 10+13 & 10+13 \\
\hline
   \end{tabular}
   \caption{Monte Carlo samples at $\sqrt{s} = 13\tev$ for 2017+2018
     setups combined. Samples are loosely sorted by category, each
     sample within one category is produced with the headlined
     generator setup. For each sample are given: a short description
     of the process, the ATLAS MC data set number (DSID), the used
     value of the generator cross section times any branching and
     filter efficiencies
     ($\sigma{\cdot} \text{BR}{\cdot}\epsilon_\mathrm{filter}$) with
     the theoretical uncertainty in percent (``th. unc.''), and
     finally the number of events analysed after skimming at
     derivation production ($N^\mathrm{skim}_\mathrm{evt}$) as well as
     the number of events originally processed and simulated
     ($N^\mathrm{unskim}_\mathrm{evt}$). Event numbers are given as sum, where the first number relates to the 2017 conditions and the second number to the 2018 conditions. The MC equivalent luminosity
     $N^\mathrm{unskim}_\mathrm{evt}/(\sigma{\cdot}\text{BR}{\cdot}\epsilon_\mathrm{filter})$
     is generally in excess of $10-20$ times the data luminosity.}
    \label{tab:samples13}
  \end{center}
\end{table}

\begin{table}[htbp]
  \small
  \begin{center}
    \begin{tabular}{l|l|l|l|l}
      \hline
      Process & Data set & $\sigma{\cdot}
      \text{BR}{\cdot}\epsilon_\mathrm{filter}$ [nb] (th. unc.)
      & $N^\mathrm{skim}_\mathrm{evt}\,[10^6]$
      & $N^\mathrm{unskim}_\mathrm{evt}\,[10^6]$\\
      \hline
      \multicolumn{5}{l}{\textbf{\Powheg+\Pythia NLO $W,Z$ samples}}\\
      \hline
$ W^{+} \to e^{+}\nu $ & 603391 & 4.220 (5\%) & 12 & 12 \\
$ W^{+} \to \mu^{+}\nu $ & 603392 & 4.220 (5\%) & 12 & 12 \\
$ W^{+} \to \tau^{+}\nu $ & 603393 & 4.220 (5\%) & 6 & 6 \\
$ W^{-} \to e^{-}\bar{\nu} $ & 603394 & 2.756 (5\%) & 8 & 8 \\
$ W^{-} \to \mu^{-}\bar{\nu} $ & 603395 & 2.756 (5\%) & 8 & 8 \\
$ W^{-} \to \tau^{-}\bar{\nu} $ & 603396 & 2.756 (5\%) & 4 & 4 \\\hline
$ Z \to ee $ $m>60\GeV$ & 603397 & 0.644 (5\%) & 4 & 4 \\
$ Z \to \mu\mu $ $m>60\GeV$ & 603398 & 0.644 (5\%) & 4 & 4 \\
$ Z \to \tau\tau $ $m>60\GeV$ & 603399 & 0.644 (5\%) & 1 & 1 \\
$ Z \to ee $ $20<m<60\GeV$ & 603400 & 0.522 (5\%) & 0.7 & 0.7 \\
$ Z \to \mu\mu $ $20<m<60\GeV$ & 603401 & 0.522 (5\%) & 0.7 & 0.7 \\
      \hline
      \multicolumn{5}{l}{\textbf{\Powheg+\Pythia MiNNLOPS $W,Z$ samples}}\\
      \hline
$ W^{+} \to e^{+}\nu $ & 604351 & 4.342 (5\%) & 11.3 & 11.3 \\
$ W^{+} \to \mu^{+}\nu $ & 604352 & 4.342 (5\%) & 11.3 & 11.3 \\
$ W^{+} \to \tau^{+}\nu $ & 604353 & 4.342 (5\%) & 5.6 & 5.4 \\
$ W^{-} \to e^{-}\bar{\nu} $ & 604348 & 2.876 (5\%) & 7.5 & 7.5 \\
$ W^{-} \to \mu^{-}\bar{\nu} $ & 604349 & 2.876 (5\%) & 7.5 & 7.5 \\
$ W^{-} \to \tau^{-}\bar{\nu} $ & 604350 & 2.876 (5\%) & 3.6 & 3.4 \\\hline
$ Z \to ee $ $m>60\GeV$ & 604354 & 0.654 (5\%) & 3.8 & 3.8 \\
$ Z \to \mu\mu $ $m>60\GeV$ & 604355 & 0.654 (5\%) & 3.8 & 3.8 \\
$ Z \to \tau\tau $ $m>60\GeV$ & 604356 & 0.654 (5\%) & 0.94 & 0.94 \\
      \hline
      \multicolumn{5}{l}{\textbf{\Powheg+\Pythia Diboson samples}}\\
      \hline
$ WW (lvlv) $ & 604493 & 0.00265 (10\%)  & 0.024 & 0.03 \\
$ WW (lvqq) $ & 604494 & 0.01102 (10\%)  & 0.057 & 0.1 \\
$ WZ (lvll) $ & 604495 & 0.00032 (10\%)  & 0.009 & 0.01 \\
$ WZ (lvqq) $ & 604496 & 0.00232 (10\%)  & 0.017 & 0.03 \\
$ WZ (lvvv) $ & 604497 & 0.00063 (10\%)  & 0.005 & 0.01 \\
$ WZ (qqll) $ & 604498 & 0.00195 (10\%)  & 0.007 & 0.02 \\
$ ZZ (llll) $ & 604499 & 0.00022 (10\%)  & 0.004 & 0.01 \\
$ ZZ (qqll) $ & 604500 & 0.00195 (10\%)  & 0.007 & 0.02 \\
$ ZZ (vvll) $ & 604501 & 0.00025 (10\%)  & 0.005 & 0.01 \\
      \hline
      \multicolumn{5}{l}{\textbf{\Powheg+\Pythia top samples}}\\
      \hline
$ t\bar{t} $ & 410470 & 0.0689 $\times$ 0.544 (7\%)  & 1.9 & 2.8 \\
$ t $ (s-chan.) & 410644 & 0.00054 (10\%)  & 0.006 & 0.01 \\
$ \bar{t} $ (s-chan.) & 410645 & 0.00028 (10\%) & 0.006 & 0.01 \\
$ t $ (t-chan.) & 410658 & 0.00545 (10\%)  & 0.6 & 1.0 \\
$ \bar{t} $ (t-chan.) & 410659 & 0.00271 (10\%)  & 0.6 & 1.0 \\
$ W\bar{t} $ & 410647 & 0.0030(10\%)  & 0.012 & 0.03 \\
$ Wt $ & 410646 & 0.0030 (10\%)  & 0.012 & 0.03\\
      \hline
      \multicolumn{5}{l}{\textbf{\Pythia $J/\psi \to \mu\mu$ samples}}\\
      \hline
Prompt & 801912 & -- & 8 & 8 \\
Non-prompt & 801917 & -- & 8 & 8\\
\hline
   \end{tabular}
   \caption{Monte Carlo samples at $\sqrt{s} = 5\tev$. Samples are loosely sorted by category, each
     sample within one category is produced with the headlined
     generator setup. For each sample are given: a short description
     of the process, the ATLAS MC data set number (DSID), the generator
     cross section times any branching and
     filter efficiencies
     ($\sigma{\cdot} \text{BR}{\cdot}\epsilon_\mathrm{filter}$) with
     the theoretical uncertainty in percent (``th. unc.''), and
     finally the number of events analysed after skimming at
     derivation production ($N^\mathrm{skim}_\mathrm{evt}$) as well as
     the number of events originally processed and simulated
     ($N^\mathrm{unskim}_\mathrm{evt}$). The MC equivalent luminosity
     $N^\mathrm{unskim}_\mathrm{evt}/(\sigma{\cdot}\text{BR}{\cdot}\epsilon_\mathrm{filter})$
     is generally in excess of $10-20$ times the data luminosity.}
    \label{tab:samples5}
  \end{center}
\end{table}
\FloatBarrier

\section{Data and MC Processing}
\label{sec:processing}

All data and MC samples are processed with the release 21 ATLAS
software. Data is processed through STDM4/6 as described above.  MC
samples are processed through STDM4, with unskimmed processing for all
signal $W$ and $Z$ samples.

All samples are further on processed to a common \texttt{MiniTree}
format using\\ \href{https://gitlab.cern.ch/atlas-physics/sm/wz/mw-lowmu/MiniTreeMaker13TeV}{https://gitlab.cern.ch/atlas-physics/sm/wz/mw-lowmu/MiniTreeMaker13TeV}.

The \texttt{MiniTree} outputs are stored in the grid dataset
\texttt{user.fballi.test.ds\_December2024\_b}
and the most current processing version of all samples is
\texttt{December2024}.

Final sample processing including application of calibrations and
corrections is done in the common \texttt{HistMaker} software,
\href{https://gitlab.cern.ch/atlas-physics/sm/wz/mw-lowmu/HistMaker}{https://gitlab.cern.ch/atlas-physics/sm/wz/mw-lowmu/HistMaker}.

\section{$z_\mathrm{vtx}$ reweighting}
\label{sec:zvtx}

The $z$-position of the primary vertex, which was selected as the hard interaction (highest $\sum
\pt^2$ of all associated tracks) was compared between data and MC for \Zee and
\Zmm selections to ensure a well-understood sample composition. 
For the 5 \TeV{} sets the simulation is specifically adapted and
perfectly matched to the data out-of-the box. 
The 13 \TeV{} set data and simulation are not well matched (very similar to
the standard Run 2 high-$\mu$ set) and
the 2017 and 2018 13 \TeV{} data sets feature slightly different beam
spots. The 13 \TeV{} MC is reweighted using the data/MC ratio from
the displayed \Zee and \Zmm selections, which by construction makes
the two distributions match well. More details can be found in the \ptw supporting note~\cite{Kretzschmar:2657141}.

\todo{Is there an impact of this correction on the mW?}

\section{Calibration of vertex selection efficiency in \Wlnu events}

Primary vertices are reconstructed requiring at least two
charged-particle tracks~\cite{ATL-PHYS-PUB-2015-026} and the hard-scatter
vertex is chosen as the one with the largest sum of squared track transverse
momenta. Lepton candidates are required to originate from this
vertex. The significance of the track's transverse impact parameter
calculated relative to the beam line, $|d_0/\sigma_{d_{0}}|$, must be
smaller than three for muons and smaller than five for
electrons. Furthermore, the longitudinal impact parameter, $z_0$ (the
difference between the $z$-coordinate of the point on the track at
which $d_0$ is defined and the longitudinal position of the primary
vertex), is required to satisfy $|z_0 \cdot \sin(\theta)| <$ 0.5 mm.

The efficiency of this requirement in \Zll events is generally above $99\%$
and included in the tag-and-probe procedures described
above. However, for \Wboson-boson events at low \pT, the
reconstruction of the correct primary vertex is more challenging as it
requires at least one additional charged-particle track from hadronic
activity in addition to the lepton. A dedicated measurement of this
efficiency is therefore derived for the \ptw analysis and is briefly described in Ref.~\cite{Atmani:2632159}. It is performed in
\Zll events as function of the number of reconstructed additional tracks, i.e. excluding the leptons, and the boson \pT. The data measurement can
then be applied to the \Wln simulation. The vertex reconstruction inefficiency in \Wboson-boson events with \ptw $< 5\,\GeV$ at
$\sqrt{s} = 5\,\TeV\,(13\,\TeV)$ is found to be about 5\% (3\%) for
the electron channels and 3\% (2\%) for the muon channels. It decreases towards higher \ptw,  and asymptotically reaches about $1\%$ in the
electron channel and $0.5\%$ in the muon channel. The inefficiency is larger in the electron channels because bremsstrahlung complicates the track reconstruction and track-to-vertex association.
The modelling of the soft activity at low \Wboson\ boson \pt\ is the main
source of uncertainty on the vertex efficiency and its effect is estimated by observing
the closure between different generators.

\todo{Check that this generator uncertainty is sufficient for mW analysis}

